\section{Hard artificial life: embodiment, adaptation, and collective construction}
\label{sec:hard-life}

Hard artificial life places synthetic organization in physical machines, where
structure, control, contact, and available power jointly determine behavior.
Its relevance to a predominantly soft-ALife review is methodological: embodiment
tests which modeled mechanisms survive material implementation. Three primary
studies illustrate distinct achievements---evolutionary design, adaptation within
an individual's lifetime, and distributed construction. They do not collectively
establish an autonomous population of reproducing robots. Keeping these
achievements separate makes their contributions easier to compare with digital
evolution.

\subsection{Joint evolution of morphology and control}

\citet{hardlipson2000} evolved robotic bodies and neural controllers together in
simulation. Mutation added, removed, or modified structural and control
components; selection rewarded displacement over a prescribed interval. Selected
designs were converted into printable structures and fabricated in thermoplastic,
after which motors were inserted manually. Three physical examples produced
locomotion corresponding to their simulated mechanisms. This demonstrated a
connection between evolutionary search and physical fabrication, with morphology
and control selected as integrated units rather than treating the body as fixed.

Transfer was imperfect. The manuscript's Table~I reports substantial distance
differences for two examples, attributed to slipping and an inadequate friction
model. The \emph{reality gap} is therefore visible within this demonstration:
qualitative transfer of a mechanism need not preserve quantitative performance.
Its autonomy also has a defined boundary. Fitness evaluation and selection took
place in simulation; fabrication used supplied equipment and materials, and
assembly retained a human step. The study supports automated embodied design,
not an ecological population manufacturing its own descendants
\citep{hardlipson2000}.

\subsection{Self-modeling and adaptation after damage}

\citet{hardbongard2006} investigated a different problem: recovering useful
behavior after an existing robot's body changes. A four-legged robot alternated
physical actions with synthesis of candidate body models. Exploratory actions
were chosen to distinguish competing models; a selected model then supported
optimization of a locomotion sequence. After part of a leg was removed, revised
models enabled an alternative gait. Comparisons against two random-exploration
baselines used the same number of physical actions, and model-driven exploration
improved inference within the tested hypothesis space.

The distinction from Darwinian evolution is central. Competing internal models
and optimized actions belong to an adaptation procedure, not a lineage of
reproducing physical individuals. Furthermore, the model family was constrained:
after damage, topology remained fixed while limb dimensions could change. The
reported locomotion demonstration used a model selected from repeated modeling
experiments. These qualifications retain the achievement---active inference can
guide compensatory behavior---while avoiding a claim of unrestricted self-knowledge,
physical regeneration, or inherited evolutionary improvement
\citep{hardbongard2006}.

\subsection{Distributed rules and externally specified collective form}

\citet{hardrubenstein2014} demonstrated programmable shape formation with
1024 Kilobots. Local communication, distance estimation, edge following,
gradients, and distributed localization supported assembly despite noisy motion
and sensing. Each robot received the same algorithm and target-shape image;
a user positioned four specially programmed seed robots to establish the
assembly's origin and orientation. The full swarm formed two shapes, each in
roughly twelve hours. Additional smaller-swarm trials examined consistency and
accuracy. Reported trials completed without charging, reprogramming, or other
human intervention during execution, although final shapes showed warping and
packing defects.

This is a strong demonstration of decentralized execution under physical
variability. It is also compatible with centralized specification of the desired
outcome: local interactions implement a supplied objective. The resulting
arrangement reorganizes existing robots; it does not manufacture new robots or
transmit variable inherited programs through reproduction. The paper also treats
physical attachment into stable structures as a future direction. Consequently,
``self-assembly'' here should be read according to the demonstrated collective
shape-formation task, rather than as evidence for a self-producing machine
organism \citep{hardrubenstein2014}.

\subsection{What embodiment contributes to soft artificial life}

The following comparison is a synthesis of these cases, not an additional
experimental result. A useful transfer is to examine the coupling between an
individual's organization and the operations available to it. In soft systems,
representation, instruction semantics, and environmental access jointly delimit
what can evolve, just as body and controller jointly delimit robotic behavior.
Interventions can ask which apparent capability disappears when one coupling is
changed. Adaptive performance should also be separated from its mechanism:
within-lifetime learning, population selection, and a designer's search may all
improve an observable score through different causal processes.

Physical embodiment additionally requires an explicit accounting boundary.
Feedstock, manufacturing equipment, electronics, power supplies, and maintenance
cannot be omitted when evaluating material autonomy. A robot's operational
independence during a trial does not close its production cycle. Soft experiments
can similarly state which resources and services their environment supplies, but
a digital resource budget is not thereby a measurement of physical energy or
manufacturing feasibility. Conversely, hardware noise need not be copied into a
digital kernel merely to make it appear biological. The transferable requirement
is to justify the abstraction and test sensitivity to it.

These historical exemplars therefore contribute tests of realization,
adaptation, and collective reliability to the review's main computational focus.
They provide neither a universal recipe for artificial life nor evidence that
open-ended evolution follows from embodiment. This bounded selection does not
survey the full contemporary frontier of evolutionary robotics.
